\documentclass[aspectratio=169]{beamer}
% Shared Beamer configuration for MUS2640 lecture decks (included after \documentclass).
\usepackage[utf8]{inputenc}
\usepackage[T1]{fontenc}
\usepackage{lmodern}
\usepackage{graphicx}
\usepackage{booktabs}

\usetheme{Madrid}
\usecolortheme{dolphin}
\usefonttheme{professionalfonts}

\setbeamercovered{transparent}
\setbeamertemplate{navigation symbols}{}
\setbeamertemplate{footline}[frame number]

\newcommand{\coursename}{MUS2640 · Sensing Sound and Music}
\newcommand{\instituteuo}{University of Oslo · Department of Musicology / RITMO}


\title{Week 1 · Tuning in}
\subtitle{\coursename}
\author{}
\institute{\instituteuo}
\date{}

\begin{document}

\begin{frame}
  \titlepage
\end{frame}

\begin{frame}{Aims for this session}
  \begin{itemize}
    \item Clarify what \emph{sensing}, \emph{sound}, and \emph{music} mean in this course.
    \item Map musicology, psychology, and technology—and levels of disciplinarity.
    \item Ground the term \emph{embodied music cognition} and preview multimodality.
  \end{itemize}
\end{frame}

\begin{frame}{Roadmap (2 hours)}
  \begin{enumerate}
    \item Course goals and your backgrounds (discussion).
    \item Interdisciplinary lenses: concepts, theories, methods.
    \item Embodied cognition \& multimodality (overview).
    \item Tools and independent work for the week.
  \end{enumerate}
\end{frame}

\begin{frame}{Two guiding questions}
  \begin{itemize}
    \item How do different disciplines ask questions about sound and music?
    \item Where do bodies, technologies, and cultures meet when we ``sense'' music?
  \end{itemize}
\end{frame}

\begin{frame}{Disciplinarity (sketch)}
  \begin{itemize}
    \item Intra- vs multi- vs inter- vs trans-disciplinary collaboration.
    \item Musicology: repertoires, history, analysis, ethnography.
    \item Psychology: perception, cognition, behaviour, experiments.
    \item Technology: capture, representation, analysis, interaction.
  \end{itemize}
\end{frame}

\begin{frame}{Embodied music cognition (preview)}
  \begin{itemize}
    \item Listening is not only ``inner hearing''—movement, gesture, and context matter.
    \item \emph{Musicking}: music as activity involving many roles and spaces.
    \item Multimodality: hearing + vision + touch/proprioception (developed later).
  \end{itemize}
\end{frame}

\begin{frame}{Suggested in-class activities}
  \begin{itemize}
    \item Short glossary: define \emph{sensing}, \emph{sound}, \emph{music} in your own words (pairs $\rightarrow$ plenary).
    \item Try one tool from the course list (optional demo): Sonic Visualiser, Audacity, or an app—reflect on what it reveals vs hides.
  \end{itemize}
\end{frame}

\begin{frame}{Outside class (about 6 hours)}
  \begin{itemize}
    \item Read \texttt{tuning-in} chapter; note terms new to you.
    \item Short reflective note on how you use AI tools (if at all) for learning.
    \item Skim the course schedule table in the \textbf{Introduction} chapter for the semester rhythm.
  \end{itemize}
\end{frame}

\begin{frame}{Summary}
  \begin{itemize}
    \item This course connects \textbf{listening bodies}, \textbf{psychological evidence}, and \textbf{technical mediation}.
    \item Next week: ways of \textbf{describing and analysing listening} (soundscapes, objects, artistic experiments).
  \end{itemize}
\end{frame}

\end{document}
